\subsection{Tree Traversal}
\begin{frame}{하나의 Tree로 네 Traversal 비교}
\centering\travtree{}{}{}{}{}{}{}\\[2mm]
\begin{tabular}{ll}Preorder&A B D E C F G\\Inorder&D B E A F C G\\Postorder&D E B F G C A\\Level-order&A B C D E F G\end{tabular}
\end{frame}
\begin{frame}{Preorder: Node를 먼저 처리}
\small\begin{algorithmic}[1]
\Procedure{Preorder}{$x$}
\If{$x=\NIL$}\State\Return\EndIf
\State \Call{Visit}{$x$}
\State \Call{Preorder}{$x.left$}
\State \Call{Preorder}{$x.right$}
\EndProcedure
\end{algorithmic}
\mission{현재 node를 처리한 뒤 left, right subtree에 같은 mission을 적용한다.}
\end{frame}
\begin{frame}{Preorder Animation}
\centering
\only<1|handout:0>{\travtree{st current}{}{}{}{}{}{}}\only<2|handout:0>{\travtree{st done}{st current}{}{}{}{}{}}\only<3|handout:0>{\travtree{st done}{st done}{}{st current}{st current}{}{}}\only<4|handout:0>{\travtree{st done}{st done}{st current}{st done}{st done}{}{}}\only<0|handout:1>{\travtree{st done}{st done}{st done}{st done}{st done}{st done}{st done}}\\
\only<1|handout:0>{visited: A}\only<2|handout:0>{visited: A, B}\only<3|handout:0>{visited: A, B, D, E}\only<4|handout:1>{visited: A, B, D, E, C, F, G}
\end{frame}
\begin{frame}{Inorder: 두 Subtree 사이에서 처리}
\small\begin{algorithmic}[1]
\Procedure{Inorder}{$x$}
\If{$x=\NIL$}\State\Return\EndIf
\State \Call{Inorder}{$x.left$}
\State \Call{Visit}{$x$}
\State \Call{Inorder}{$x.right$}
\EndProcedure
\end{algorithmic}
\mission{left subtree를 마친 뒤 node, 그다음 right subtree를 처리한다.}
\end{frame}
\begin{frame}{Inorder Animation}
\centering
\only<1|handout:0>{\travtree{}{}{ }{st current}{}{}{}}\only<2|handout:0>{\travtree{}{st current}{}{st done}{st current}{}{}}\only<3|handout:0>{\travtree{st current}{st done}{}{st done}{st done}{}{}}\only<4|handout:1>{\travtree{st done}{st done}{st done}{st done}{st done}{st done}{st done}}\\
\only<1|handout:0>{visited: D}\only<2|handout:0>{visited: D, B, E}\only<3|handout:0>{visited: D, B, E, A}\only<4|handout:1>{visited: D, B, E, A, F, C, G}
\end{frame}
\begin{frame}{Postorder: Subtree를 모두 마친 뒤 처리}
\small\begin{algorithmic}[1]
\Procedure{Postorder}{$x$}
\If{$x=\NIL$}\State\Return\EndIf
\State \Call{Postorder}{$x.left$}
\State \Call{Postorder}{$x.right$}
\State \Call{Visit}{$x$}
\EndProcedure
\end{algorithmic}
\mission{두 subtree의 결과가 준비된 뒤 node를 처리한다.}
\end{frame}
\begin{frame}{Postorder Animation}
\centering
\only<1|handout:0>{\travtree{}{}{}{st current}{}{}{}}\only<2|handout:0>{\travtree{}{st current}{}{st done}{st current}{}{}}\only<3|handout:0>{\travtree{}{st done}{st current}{st done}{st done}{st done}{st done}}\only<4|handout:0>{\travtree{st current}{st done}{st done}{st done}{st done}{st done}{st done}}\only<0|handout:1>{\travtree{st done}{st done}{st done}{st done}{st done}{st done}{st done}}\\
\only<1|handout:0>{visited: D}\only<2|handout:0>{visited: D, E, B}\only<3|handout:0>{visited: D, E, B, F, G, C}\only<4|handout:1>{visited: D, E, B, F, G, C, A}
\end{frame}
\begin{frame}{DFS Traversal의 비용}
\begin{itemize}
\item 각 node에서 $\Theta(1)$ work, 정확히 한 번 방문 → time $\Theta(n)$
\item recursive stack는 현재 root-to-node path → $O(h)$
\item balanced tree: $O(\log n)$ stack; degenerate tree: $O(n)$
\end{itemize}
\end{frame}
\begin{frame}{Level-Order Pseudocode}
\small\begin{algorithmic}[1]
\Procedure{LevelOrder}{$root$}
\If{$root=\NIL$}\State\Return\EndIf
\State $Q\gets$ empty queue; \Call{Enqueue}{$Q,root$}
\While{$Q$ is not empty}
\State $u\gets\Call{Dequeue}{Q}$; \Call{Visit}{$u$}
\If{$u.left\ne\NIL$}\State\Call{Enqueue}{$Q,u.left$}\EndIf
\If{$u.right\ne\NIL$}\State\Call{Enqueue}{$Q,u.right$}\EndIf
\EndWhile
\EndProcedure
\end{algorithmic}
\end{frame}
\begin{frame}{Level-Order Queue Animation}
\centering\travtree{}{}{}{}{}{}{}\\[-1mm]
\begin{tabular}{ccll}\toprule
step&current&queue before $\to$ after&visited\\\midrule
\only<1|handout:0>{0&--&[] $\to$ [A]&[]}
\only<2|handout:0>{1&A&[A] $\to$ [B,C]&[A]}
\only<3|handout:0>{2&B&[B,C] $\to$ [C,D,E]&[A,B]}
\only<4|handout:1>{3–7&C,D,E,F,G&[C,D,E] $\to\cdots\to$ []&[A,B,C,D,E,F,G]}\\\bottomrule
\end{tabular}
\end{frame}
\begin{frame}{BFS Queue의 공간}
\begin{itemize}
\item time은 node와 edge를 한 번 다루므로 $\Theta(n)$
\item queue에는 frontier가 저장된다.
\item maximum width를 $w$라 하면 auxiliary space $O(w)$
\item perfect tree에서는 마지막 level 때문에 $w=\Theta(n)$ 가능
\end{itemize}
\end{frame}
\begin{frame}{Expression Tree}
\centering\begin{tikzpicture}[level distance=10mm,level 1/.style={sibling distance=46mm},level 2/.style={sibling distance=22mm},level 3/.style={sibling distance=11mm}]
\node[st current]{*}
 child{node[st node]{+}child{node[st node]{x}}child{node[st node]{y}}}
 child{node[st node]{/}child{node[st node]{+}child{node[st node]{a}}child{node[st node]{b}}}child{node[st node]{c}}};
\end{tikzpicture}
\end{frame}
\begin{frame}{Traversal이 Notation을 만든다}
\begin{description}
\item[Prefix / preorder] \texttt{* + x y / + a b c}
\item[Infix / inorder] \texttt{(x + y) * ((a + b) / c)}
\item[Postfix / postorder] \texttt{x y + a b + c / *}
\end{description}
\begin{alertblock}{괄호가 필요한 이유}단순 inorder \texttt{x + y * a + b / c}는 원래 grouping을 보존하지 않는다. 각 operator subtree를 괄호로 감싼다.\end{alertblock}
\end{frame}
\begin{frame}{Traversal Checkpoint}
\begin{enumerate}
\item fixed tree의 inorder에서 A 바로 앞 key는?
\item postorder의 마지막 node는?
\item recursive DFS와 BFS의 auxiliary space는?
\item expression tree의 postfix는?
\end{enumerate}
\begin{exampleblock}{답}E; root A; $O(h)$와 $O(w)$; \texttt{x y + a b + c / *}.\end{exampleblock}
\end{frame}
